\lecture{14}{Thu 26 Mar 2026 12:31}{Sheaves}

\section{Sheaf cohomology}

Let $\mathcal{F} $ be a sheaf on a space $X$ valued in a concrete abelian category. Let $U\subseteq X$ be open. Then there is a map to the product of stalks
\[
	\mathcal{F} (U)\to \prod_{x\in U} \mathcal{F} _x
\]
given by $s\mapsto \left\{ s_x \right\}_{x\in U} $, where $s_x$ is the germ of $s$ at $x$. This map is injective since its kernel is trivial: if $s_x=0$, then there is $V\subseteq U$ containing $x$ such that $s|V=0$; the gluing axiom yields $s=0$.

\begin{notation}\label{4.1.1}
	We write the category of presheaves on a space $X$ valued in $\mathcal{C} $ as $\Pshf_X(\mathcal{C} )$. We often leave the $\mathcal{C} $ implicit and write $\Pshf _X$. Similarly, we write $\Shf _X(\mathcal{C} )=\Shf _X$ for the full subcategory of $\Pshf_X(\mathcal{C} )$ spanned by sheaves.
\end{notation}


\begin{construction}\label{4.1.2}
	Let $\mathcal{F} $ be a presheaf on a space $X$. Let $U : \Shf_X  \to \Pshf_X$ be the forgetful functors. Then the \emph{sheafification} $\mathcal{F} ^+$ of $\mathcal{F} $ is the sheaf which is universal in the sense that there is a natural transformation $i : \mathcal{F}  \Rightarrow \mathcal{F} ^+$ whose precomposition induces isomorphisms
	\[
		i^* : \Shf_X (\mathcal{F} ^+,\mathcal{G} )\cong \Pshf_X (\mathcal{F} ,U(\mathcal{G} ))
	\]
	for all sheaves $\mathcal{G} $, natural in $\mathcal{F} $ and $\mathcal{G} $. More precisely, $-^+ \dashv U$, meaning the sheafification of a sheaf is the free sheaf on a presheaf.

	A concrete discussion is as follows. Let $U\subseteq X$ be open, and define $F^+(U)$ as the set of all functions $t : U\to \prod_{x\in U} \mathcal{F} _x$ having the property that for all $x\in U$, there exists $V\ni x$ and a section $s\in \mathcal{F} (U)$ such that $t(x)=s_x$. This is called the \emph{Etal\'e space}.

	Define the set
	\[
		E\coloneqq \coprod_{x\in X} \mathcal{F} _x
	.\]
	Topologize it as follows. If $s\in \mathcal{F} (U)$, then define the map $\widetilde{s} : U \to E$ by $\widetilde{s} (x)=s_x$. The topology on $E$ is generated by taking the following as a basis:
	\[
		\mathcal{B}\coloneqq \left\{ \widetilde{s} (U)\mid s\in \mathcal{F} (V),U,V\in \Open(X)  \right\} 
	.\]
	The projection $\pi : E \to X$ is then a local homeomorphism, and since it is surjective, we can consider sections $\Gamma (U,E)$ of an open set $U\subseteq X$. The sheafification $\mathcal{F} ^+$ can be defined as the sheaf of sections
	\[
		\mathcal{F}^+ \cong \Gamma (-,E)
	.\]
\end{construction}

\begin{remark}\label{4.1.3}
	Sheaves are stable under \emph{most} relevant (co)limits such that pointwise construction of (co)limits yields an abelian category structure for $\Shf _X(\mathcal{A} )$ whenever $\mathcal{A} $ is abelian. For example, if $\phi : \mathcal{F}  \to \mathcal{G} $ is a natural transformation, then
	\[
		\Ker \phi (U) \coloneqq \Ker (\phi_U : \mathcal{F} (U) \to \mathcal{G} (U))
	\]
	is a sheaf. To show this, let $\left\{ U_i \right\} _{i\in I} $ an open cover of $U$, let $\left\{ s_i \right\}_{i\in I}$ with $s_i\in \Ker \phi(U_i)=\Ker \phi _{U_i} $. Assume that $s_i|U_i \cap U_j = s_j|U_i \cap U_j$ for all $i,j$. Then we identify $\Ker \phi _{U_i} \subseteq \mathcal{F} (U_i)$ to obtain that $s_i|U_i \cap U_j = s_j|U_i \cap U_j$ as elements of $\mathcal{F} $, since the restrictions agree in the kernel. Thus there exists $s\in \mathcal{F} (U)$ with $s|U_i=s_i$. We show $\phi _U(s)=0$. Indeed, by naturality of $\phi $ we obtain
	\[
		\phi _U(s)|U_i = \phi _U(s_i) = 0
	.\]
	Thus $\phi _U(s)$ is given by gluing together $0\in \mathcal{G} (U_i)$ along the open cover $\left\{ U_i \right\} _{i\in I} $, locality yields $\phi _U(s)=0$.

	However, $\Shf _X$ is not stable under all colimits. For example, $\Im \phi $ is not a sheaf. Since $\Im = \Ker \coker $, this tells us that the cokernel need not be a sheaf. We instead define the image $\Im \phi $ to be the sheafification of the pointwise image of $\phi $.

	More generally, we note the following: if $\mathcal{F} $ is a presheaf, then it is a sheaf if and only if for all opens $U_a,U_b\subseteq X$,
	\[
		\mathcal{F} (U) \cong \mathcal{F} (U_a)\times_{\mathcal{F} (U_a \cap U_b)} \mathcal{F} (U_b)
	.\]
	Thus $\Shf _X$ is stable under limits since limits commute, but not necessarily under colimits. Thus we by convention re-define colimits in $\Shf _X$ by taking the sheafification of the relevant colimit in $\Pshf $. Since sheafification is left-adjoint, it preserves colimits, and thus is indeed a colimit.
\end{remark}

\begin{definition}\label{4.1.4}
	Since the image of a morphism of sheaves (as a categorical definition) is \emph{not} the image as a morphism of presheaves, we define a morphism $\phi : \mathcal{F}  \to \mathcal{G} $ of sheaves to be \emph{surjective/epi} if $\Im \phi \cong \mathcal{G} $.
\end{definition}
\begin{exercise}\label{4.1.5}
	With notation as in \cref{4.1.4}, the following are equivalent.
	\begin{enumerate}
		\item $\phi $ is surjective.
		\item $\phi _x : \mathcal{F} _x \to \mathcal{G} _x$ is surjective as a function for all $x\in X$.
		\item For all open $U\subseteq X$ and $s\in \mathcal{G} (U)$, there exists $x\in U$, an open set $x\in V \subseteq U$, and a section $t\in \mathcal{F} (V)$ such that $\phi (t)=s|V$. In particular, locally a surjective morphism of sheaves is surjective in the usual sense.
	\end{enumerate}
\end{exercise}

\begin{example}\label{4.1.6}
	Consider the sheaf $\mathcal{O} ^*_X$ of nowhere vanishing holomorphic functions on a complex manifold $X$. This takes values in multiplicative abelian groups, so we have a natural transformation $\exp : \mathcal{O} _X \to \mathcal{O} ^*_X$ which maps $f : U\subseteq X \to \mathbb{C} $ to $\exp (2\pi if) : U \to \mathbb{C} $. This map is very important. We note that $f$ is in the kernel of $\exp _U$ if and only if $f(z)=n$ for some $n\in\mathbb{Z} $ and all $z\in U$. Hence $\Ker \exp \cong \mathbb{Z} _X$ is the constant sheaf on $X$.

	We can also see that $\exp $ is surjective. More precisely, for all $x\in X$ and $g\in \mathcal{O} _X^*$ defined near $x$, we can find some $f\in \mathcal{O} _X^*$ defined in a neighborhood of $x$ with $\exp(2\pi if)=g$. Hence we are claiming that locally a logarithm can be defined.

	To see this, let $g(x)\in\mathbb{C} $. Since this is nonzero, there is a neighborhood which does not contain $0$, and thus sufficiently near as to avoid a branch cut. Thus $\ln g$ can be defined on the preimage $g^{-1} (U)$ of some sufficiently small open neighborhood $U\ni g(x)$.

	It follows that
	\[\begin{tikzcd}[row sep = 2.5em, column sep = 2.5em]
	0\ar[r]  & \mathbb{Z} _X\ar[r]  & \mathcal{O} _X\ar[" \exp  ", r]  & \mathcal{O} _X^*\ar[r]  & 0
	\end{tikzcd}\]
	is an exact sequence of sheaves.
\end{example}

\begin{example}\label{4.1.7}
	Let $X$ be a complex manifold. Recall that $\mathcal{A} ^{p,q}_X$ is the sheaf of smooth $(p,q)$-forms on $X$. Write $\mathcal{A} ^{p,\bullet} $ for the complex
	\[\begin{tikzcd}[row sep = 2.5em, column sep = 2.5em]
	0\ar[r]  & \mathcal{A} ^{p,0}_X \ar[" \overline{\partial }  ", r]  & \mathcal{A} ^{p,1}_X \ar[" \overline{\partial }  ", r]  & \mathcal{A} ^{p,2}_X \ar[" \overline{\partial }  ", r]  & \cdots .
	\end{tikzcd}\]
	As seen before, this is a complex since $\overline{\partial } ^2=0$. We claim this sequence is exact at $\mathcal{A} ^{p,q} $ for all $q\geq 1$. This follows from the $\overline{\partial } $-Poincare lemma, which says that on a polydisk, every $\overline{\partial } $-closed $(p,q)$ form with $q>0$ is $\overline{\partial } $-exact.
\end{example}

\begin{definition}\label{4.1.8}
	A sheaf $\mathcal{F} $ on a complex manifold is an $\mathcal{O} _X$-module if $F(U)$ is an $\mathcal{O} _X(U)$-module for all open sets $U$. More generally, if $(X,\mathcal{O} _X)$ is a locally ringed space (meaning $\mathcal{O} _X$ is a sheaf of rings, called the structure sheaf, and usually some kind of sheaf of functions), then an $\mathcal{O} _X$-module is a sheaf $\mathcal{F} $ on $X$ such that $\mathcal{F} (U)$ is an $\mathcal{O} _X(U)$-module for all open sets $U\subseteq X$.
\end{definition}

\begin{example}\label{4.1.9}
	For example, if $\pi : E \to X$ is a holomorphic vector bundle, then the sheaf of sections $\mathcal{E} =\Gamma (-,E)$ is an $\mathcal{O} _X$-module (since we can just multiply two holomorphic functions).
\end{example}

\begin{example}\label{4.1.10}
	Let $x\in X$. Define the \emph{skyscraper sheaf at $x$} as the sheaf $S_x(U)$ to be $\mathbb{C} $ if $x\in U$ and $0$ otherwise.
\end{example}

Consider a category of sheaves $\Shf _X(\mathcal{A} )$ on an abelian category. The global sections functor $\mathcal{F} \mapsto \Gamma (X,\mathcal{F} )$ is not exact (it is right adjoint to the constant sheaf functor). More concretely, surjectivity at the level of sheaves and at the level of sets of sections are \emph{different}, meaning we only get that $\Gamma (X,-)$ preserves kernels, but not surjectivity. Thus it has right derived functors.

\begin{definition}\label{4.1.11}
	Let $\mathcal{F} $ be a sheaf on $X$ valued in an abelian category $\mathcal{A} $. The \emph{sheaf cohomology} is given by the right derived functors of the global sections functor
	\[
		H^{\bullet} (X,\mathcal{F} )\coloneqq R^n\Gamma (X,\mathcal{F} )
	.\]
\end{definition}

